\documentclass[11pt]{article}

\usepackage[utf8]{inputenc}
\usepackage[T1]{fontenc}
\usepackage[spanish,es-nodecimaldot]{babel}
\usepackage{amsmath,amssymb,mathtools}
\usepackage{geometry}
\usepackage{booktabs}
\usepackage{array}
\usepackage{enumitem}
\usepackage{hyperref}
\usepackage{xcolor}
\usepackage{siunitx}
\usepackage{graphicx}
\usepackage{float}

\geometry{margin=2.5cm}
\hypersetup{
    colorlinks=true,
    linkcolor=blue!50!black,
    urlcolor=blue!50!black
}

\title{Plataforma Antivibratoria Activa para un Sensor de Ultrasonido\\Documento Ajustado al Enunciado}
\author{Base para la entrega en Overleaf}
\date{21 de abril de 2026}

\begin{document}

\maketitle
\tableofcontents
\newpage

\section{Introduccion}

En el laboratorio se realizan pruebas sobre papas usando un sensor de ultrasonido. La medicion se ve afectada por vibraciones de la mesa o base donde esta montado el sensor. Si la base vibra, la posicion del sensor cambia y la distancia medida deja de reflejar unicamente la geometria real del objeto.

Por esta razon, el problema se plantea como un problema de \textbf{rechazo de perturbaciones}. El objetivo del sistema de control es mantener lo mas inmovil posible la plataforma del sensor aunque la base vibre.

El sistema propuesto es SISO y su planta nominal es de \textbf{tercer orden}. La variable manipulada es el voltaje aplicado al actuador, mientras que la salida de interes es la posicion de la plataforma donde esta montado el sensor.

\section{1. Modelado realista del sistema}

\subsection{Descripcion fisica}

El sistema realista esta compuesto por:

\begin{itemize}[leftmargin=1.2cm]
\item una base inferior sometida a vibraciones externas;
\item una plataforma superior donde esta montado el sensor de ultrasonido;
\item un resorte entre base y plataforma;
\item un amortiguador viscoso entre base y plataforma;
\item un actuador electromecanico que aplica una fuerza de correccion.
\end{itemize}

\subsection{Supuestos del modelo realista}

\begin{itemize}[leftmargin=1.2cm]
\item la plataforma se modela como una masa concentrada;
\item el movimiento dominante se estudia en una sola direccion;
\item el resorte trabaja de forma lineal en el rango principal de operacion;
\item el amortiguamiento se aproxima como viscoso lineal;
\item el actuador se modela como un elemento electromecanico lineal;
\item el modelo realista puede incluir saturacion del actuador y restricciones fisicas adicionales en Simulink/Simscape.
\end{itemize}

\subsection{Variables principales}

\begin{center}
\small
\begin{tabular}{>{\raggedright\arraybackslash}p{2.8cm} >{\raggedright\arraybackslash}p{9.3cm}}
\toprule
\textbf{Simbolo} & \textbf{Significado} \\
\midrule
$x(t)$ & posicion absoluta de la plataforma del sensor \\
$z(t)$ & posicion absoluta de la base \\
$q(t)=x(t)-z(t)$ & desplazamiento relativo entre plataforma y base \\
$u(t)$ & voltaje aplicado al actuador \\
$i(t)$ & corriente del actuador \\
$F_a(t)$ & fuerza generada por el actuador \\
$y(t)=x(t)$ & salida regulada del sistema \\
\bottomrule
\end{tabular}
\end{center}

\subsection{Espacio para la figura del modelo realista}

\begin{figure}[H]
\centering
\includegraphics[width=0.82\textwidth]{fig_modelo_realista_simscape.png}
\caption{Modelo realista implementado en Simulink/Simscape.}
\label{fig:modelo_realista}
\end{figure}

\noindent\textbf{Texto sugerido:} En la Fig.~\ref{fig:modelo_realista} se muestra el modelo realista construido en Simulink/Simscape. En esta figura se deben identificar la base, la plataforma, el resorte, el amortiguador, el actuador y el punto donde se monta el sensor.

\section{2. Obtencion del modelo matematico}

\subsection{Modelo fisico completo}

Se define la coordenada relativa:
\begin{equation}
q(t)=x(t)-z(t).
\end{equation}

Aplicando la segunda ley de Newton:
\begin{equation}
m\ddot x = -k(x-z)-c(\dot x-\dot z)+F_a.
\end{equation}

Pasando a coordenada relativa:
\begin{equation}
m\ddot q + c\dot q + kq = F_a - m\ddot z.
\label{eq:mech_full}
\end{equation}

Para el actuador se usa:
\begin{equation}
F_a=K_f i,
\label{eq:fa}
\end{equation}
\begin{equation}
L\dot i + Ri + K_e\dot q = u.
\label{eq:elec}
\end{equation}

Sustituyendo \eqref{eq:fa} en \eqref{eq:mech_full}:
\begin{equation}
m\ddot q + c\dot q + kq = K_f i - m\ddot z,
\label{eq:full1}
\end{equation}
\begin{equation}
L\dot i + Ri + K_e\dot q = u.
\label{eq:full2}
\end{equation}

\subsection{Planta nominal para analisis y control}

Para obtener la planta nominal se anula temporalmente la perturbacion:
\[
z(t)=0,\qquad \dot z(t)=0,\qquad \ddot z(t)=0.
\]

Entonces:
\begin{equation}
m\ddot q + c\dot q + kq = K_f i,
\end{equation}
\begin{equation}
L\dot i + Ri + K_e\dot q = u.
\end{equation}

Se definen los estados:
\[
x_1=q,\qquad x_2=\dot q,\qquad x_3=i.
\]

La representacion nominal en espacio de estados queda:
\begin{equation}
\dot x = Ax+Bu,
\end{equation}
\begin{equation}
y = Cx+Du,
\end{equation}
con
\begin{equation}
A=
\begin{bmatrix}
0 & 1 & 0\\
-\dfrac{k}{m} & -\dfrac{c}{m} & \dfrac{K_f}{m}\\
0 & -\dfrac{K_e}{L} & -\dfrac{R}{L}
\end{bmatrix},
\qquad
B=
\begin{bmatrix}
0\\
0\\
\dfrac{1}{L}
\end{bmatrix},
\end{equation}
\begin{equation}
C=
\begin{bmatrix}
1 & 0 & 0
\end{bmatrix},
\qquad
D=0.
\end{equation}

La funcion de transferencia nominal se obtiene con:
\begin{equation}
G(s)=C(sI-A)^{-1}B+D.
\end{equation}

Al desarrollar:
\begin{equation}
G(s)=\frac{Q(s)}{U(s)}
=
\frac{K_f}
{(Ls+R)(ms^2+cs+k)+K_fK_e s}.
\label{eq:tf3}
\end{equation}

Expandiendo:
\begin{equation}
G(s)=
\frac{K_f}
{Lm\,s^3 + (Lc+Rm)s^2 + (Lk+Rc+K_fK_e)s + Rk}.
\end{equation}

\subsection{Analisis de estabilidad}

La ecuacion caracteristica asociada a la planta nominal es:
\[
Lm\,s^3 + (Lc+Rm)s^2 + (Lk+Rc+K_fK_e)s + Rk = 0.
\]

Definiendo:
\[
a_3=Lm,\qquad a_2=Lc+Rm,\qquad a_1=Lk+Rc+K_fK_e,\qquad a_0=Rk,
\]
el criterio de Routh-Hurwitz para un polinomio cubico exige:
\begin{equation}
a_3>0,\quad a_2>0,\quad a_1>0,\quad a_0>0,\quad a_2a_1>a_3a_0.
\end{equation}

Como los parametros fisicos son positivos, la condicion relevante es:
\begin{equation}
(Lc+Rm)(Lk+Rc+K_fK_e)>LmRk.
\end{equation}

Por tanto, la planta nominal es de \textbf{tercer orden} y su estabilidad se analiza sobre ese modelo.

\section{3. Comparacion entre modelos}

La comparacion requerida por el enunciado se hace entre:

\begin{itemize}[leftmargin=1.2cm]
\item el modelo realista implementado en Simulink/Simscape;
\item el modelo matematico aproximado usado para analisis y control.
\end{itemize}

La comparacion en lazo abierto se realiza frente a una entrada escalon y permite identificar diferencias entre:

\begin{itemize}[leftmargin=1.2cm]
\item rapidez de respuesta;
\item amortiguamiento;
\item sobreimpulso;
\item efectos fisicos no capturados por el modelo teorico.
\end{itemize}

\subsection{Espacio para grafica: lazo abierto teorico}

\begin{figure}[H]
\centering
\includegraphics[width=0.82\textwidth]{fig_lazo_abierto_tf.png}
\caption{Respuesta en lazo abierto del modelo teorico obtenido a partir de la funcion de transferencia.}
\label{fig:abierto_tf}
\end{figure}

\subsection{Espacio para grafica: lazo abierto Simscape}

\begin{figure}[H]
\centering
\includegraphics[width=0.82\textwidth]{fig_lazo_abierto_simscape.png}
\caption{Respuesta en lazo abierto del modelo realista en Simscape.}
\label{fig:abierto_simscape}
\end{figure}

\subsection{Texto para completar}

\noindent\textbf{Texto sugerido:} Las Figs.~\ref{fig:abierto_tf} y \ref{fig:abierto_simscape} permiten comparar el modelo aproximado con el modelo realista. Aqui se debe discutir en que medida la planta teorica reproduce la dinamica dominante del sistema fisico y cuales diferencias aparecen por saturaciones, restricciones u otros efectos del modelo realista.

\section{4. Diseno del controlador}

\subsection{Modelo reducido para sintonia IMC-PID}

Para la sintonia del PID se usa un modelo reducido de segundo orden:
\begin{equation}
G_r(s)=\frac{1}{ms^2+cs+k}.
\end{equation}

Con los valores dados:
\[
m=1,\qquad c=5,\qquad k=100,
\]
queda:
\begin{equation}
G_r(s)=\frac{1}{s^2+5s+100}.
\label{eq:gred}
\end{equation}

Se escoge como modelo deseado:
\begin{equation}
T_d(s)=\frac{1}{\alpha s+1}.
\end{equation}

Entonces:
\begin{equation}
C(s)=\frac{T_d(s)}{G_r(s)\left[1-T_d(s)\right]}
=
\frac{s^2+5s+100}{\alpha s}.
\label{eq:imc}
\end{equation}

\subsection{Controlador PID ideal}

La forma ideal del PID es:
\begin{equation}
C_{PID}(s)=K_c\left(1+\frac{1}{\tau_I s}+\tau_D s\right).
\end{equation}

Al comparar con \eqref{eq:imc}, se obtiene:
\begin{equation}
K_c=\frac{5}{\alpha},\qquad
\tau_I=0.05\ \text{s},\qquad
\tau_D=0.2\ \text{s}.
\end{equation}

Si se desea mantener la ganancia proporcional en:
\[
P=K_c=50,
\]
entonces:
\begin{equation}
50=\frac{5}{\alpha}
\quad\Rightarrow\quad
\alpha=0.1\ \text{s}.
\end{equation}

Por tanto, el controlador PID ideal final es:
\begin{equation}
\boxed{
C(s)=50\left(1+\frac{1}{0.05\,s}+0.2\,s\right)
}
\label{eq:pidfinal}
\end{equation}

En forma paralela:
\[
K_p=50,\qquad K_i=1000,\qquad K_d=10.
\]

\subsection{Equivalencia con el bloque PID}

Si se usa el bloque:
\begin{equation}
C_b(s)=P\left(1+\frac{1}{Is}+D\frac{N}{1+N/s}\right),
\end{equation}
la equivalencia es:
\[
P=50,\qquad I=0.05,\qquad D=0.2.
\]

El termino derivativo del bloque es filtrado:
\[
D\frac{N}{1+N/s}
=
D\frac{Ns}{s+N},
\]
por lo que para un valor inicial razonable se puede tomar:
\[
N=10.
\]

Entonces, para Simulink:
\begin{equation}
\boxed{
P=50,\qquad I=0.05,\qquad D=0.2,\qquad N=10
}
\label{eq:pidsimulink}
\end{equation}

\section{5. Evaluacion del desempeno}

La evaluacion del desempeno del controlador debe hacerse tanto sobre el modelo aproximado como sobre el modelo realista. En particular, el enunciado pide estudiar:

\begin{itemize}[leftmargin=1.2cm]
\item rechazo a perturbaciones;
\item seguimiento de una referencia escalon variable;
\item tiempo de establecimiento;
\item tiempo de subida;
\item sobreimpulso;
\item error en estado estacionario.
\end{itemize}

\subsection{Espacio para grafica: lazo cerrado Simscape}

\begin{figure}[H]
\centering
\includegraphics[width=0.82\textwidth]{fig_lazo_cerrado_simscape.png}
\caption{Respuesta en lazo cerrado del modelo realista en Simscape con el controlador IMC-PID.}
\label{fig:cerrado_simscape}
\end{figure}

\subsection{Espacio para tabla de metricas}

\begin{table}[H]
\centering
\small
\begin{tabular}{lcccc}
\toprule
\textbf{Caso} & \textbf{$t_r$} & \textbf{$t_s$} & \textbf{$M_p$} & \textbf{$e_{ss}$} \\
\midrule
Lazo abierto teorico & -- & -- & -- & -- \\
Lazo abierto Simscape & -- & -- & -- & -- \\
Lazo cerrado Simscape & -- & -- & -- & -- \\
\bottomrule
\end{tabular}
\caption{Metricas de desempeno a completar con los resultados experimentales o de simulacion.}
\label{tab:metricas}
\end{table}

\subsection{Texto para completar}

\noindent\textbf{Texto sugerido:} En la Fig.~\ref{fig:cerrado_simscape} se presenta la respuesta del sistema realista en lazo cerrado. Aqui se debe analizar si el controlador reduce la transmision de vibracion hacia la plataforma y si mantiene el sensor suficientemente estable para realizar la medicion. Los resultados cuantitativos deben resumirse en la Tabla~\ref{tab:metricas}.

\section{Espacio para bono}

El enunciado indica que se puede obtener una bonificacion adicional si el proyecto incluye una visualizacion tridimensional del sistema. Por esta razon se deja esta seccion opcional para documentar ese desarrollo.

\subsection{Visualizacion 3D opcional}

\begin{figure}[H]
\centering
\includegraphics[width=0.82\textwidth]{fig_bono_visualizacion_3d.png}
\caption{Visualizacion tridimensional opcional del sistema.}
\label{fig:bono_3d}
\end{figure}

\noindent\textbf{Texto sugerido:} En la Fig.~\ref{fig:bono_3d} se muestra la visualizacion tridimensional del sistema. Esta parte es opcional y se incluye solo si se desarrolla la bonificacion correspondiente.

\section{Conclusiones}

El proyecto cumple con lo solicitado en el enunciado: se propone un sistema SISO de al menos tercer orden, se construye un modelo realista, se deriva un modelo matematico nominal, se dise\~na un controlador por IMC-PID y se deja estructurada la comparacion entre modelos y la evaluacion del desempeno.

Con los parametros $m=1$, $c=5$ y $k=100$, y manteniendo $P=50$, el controlador final queda:
\[
C(s)=50\left(1+\frac{1}{0.05\,s}+0.2\,s\right),
\]
y su forma equivalente para el bloque PID es:
\[
P=50,\qquad I=0.05,\qquad D=0.2,\qquad N=10.
\]

\end{document}
